%% LyX 2.4.4 created this file.  For more info, see https://www.lyx.org/.
%% Do not edit unless you really know what you are doing.
\documentclass[english]{article}
\usepackage[T1]{fontenc}
\usepackage[latin9]{inputenc}
\usepackage{units}
\usepackage{enumitem}
\usepackage{amsmath}
\usepackage{amsthm}
\usepackage{amssymb}
\usepackage{cancel}
\usepackage{geometry}
\geometry{verbose,tmargin=1in,bmargin=1in,lmargin=1in,rmargin=1in}

\makeatletter

%%%%%%%%%%%%%%%%%%%%%%%%%%%%%% LyX specific LaTeX commands.
%% Because html converters don't know tabularnewline
\providecommand{\tabularnewline}{\\}

%%%%%%%%%%%%%%%%%%%%%%%%%%%%%% Textclass specific LaTeX commands.
\numberwithin{equation}{section}
\numberwithin{figure}{section}
\newlength{\lyxlabelwidth}      % auxiliary length 

%%%%%%%%%%%%%%%%%%%%%%%%%%%%%% User specified LaTeX commands.
\usepackage{fancyhdr}
\usepackage{lastpage}
\pagestyle{fancy}
\usepackage{enumitem}
\usepackage{ifthen}
%sets math fonts to be more readable
\everymath=\expandafter{\the\everymath\displaystyle}
%\DeclareMathSizes{10}{10}{10}{10}
\usepackage{multicol}

\usepackage{tikz}
\usetikzlibrary{plotmarks}
\usepackage{pgfplots}
\pgfplotsset{compat=newest}

\usepackage{calc}
\usepackage[nomessages]{fp}

\usepackage{advdate}    % Advancing/saving dates
\usepackage{datetime}   % Dates formatting
\usepackage{datenumber} % Counters for dates
% Added by lyx2lyx
\setlength{\parskip}{\smallskipamount}
\setlength{\parindent}{0pt}

\makeatother

\usepackage{babel}
\begin{document}
\renewcommand{\author}{Dr.\,James Ashe}

\newcommand{\classNum}{132}

\newcommand{\classSec}{C and K}

\newcommand{\nameType}{Name:}

\newcommand{\examType}{Midterm Solutions}

\renewcommand{\date}{\formatdate{4}{3}{2014}}

%%First page's header%%%%%%%%%%%%%%%%%%%%%%
\fancypagestyle{firststyle} %%header settings for first pagr
{
	\fancyhead[L]{MTH \classNum{}(\classSec{}) \author{}\\
	\textbf{\examType{}} \date{}}
	\fancyhead[C]{\textbf{\nameType}}
	\setlength{\headsep}{14pt}
}
%%%%%%%%%%%%%%%%%%%%%%%%%%%%%%%%%%%%%%%%%%

%%Next page's header%%%%%%%%%%%%%%%%%%%%%%%
\setlist{nolistsep}
\lhead{MTH \classNum{}(\classSec{})} 
\chead{\textbf{\nameType}} 
\cfoot{\thepage{}~of~\pageref{LastPage}} 
\setlength{\headsep}{5pt} 
%%%%%%%%%%%%%%%%%%%%%%%%%%%%%%%%%%%%%%%%%%%

%%Various settings; see the preamble for other settings%%%%%
%%\setenumerate[0]{leftmargin=12pt,itemindent=0pt,labelwidth=0pt}
\renewcommand{\labelenumi}{\textbf{\arabic{enumi}.}}
\renewcommand{\labelenumii}{\textbf{\alph{enumii}.}}
\thispagestyle{firststyle}
%%%%%%%%%%%%%%%%%%%%%%%%%%%%%%%%%%%%%%%%%%
\newcommand{\xmax}{10}
\newcommand{\xmin}{-10}
\newcommand{\xstep}{0} %must be xmin+n
\newcommand{\ymax}{10}
\newcommand{\ymin}{-10}
\newcommand{\ystep}{0} %must be ymin+n

\newcommand{\xspace}{\vspace{1.5cm}}

\emph{You must show work to receive credit. Do not use a calculator/computer
for problems with a }$\bigstar$\emph{. Each problem is of equal value. }

\textbf{Find the slope of the line passing through the given pair
of points .}
\begin{enumerate}[resume]
\item $(1,2)$ and $(3,-4)$

\begin{enumerate}[resume]
\item []$m=\frac{y_{2}-y_{1}}{x_{2}-x_{1}}=\frac{2-(-4)}{1-3}=\frac{6}{-2}=-3$\vfill
\end{enumerate}
\end{enumerate}
\textbf{Use the given conditions to write an equation for the line
in slope-intercept form, if possible.}
\begin{enumerate}[resume]
\item Slope$=-2$, passing through $(13,4)$.

\begin{enumerate}[resume]
\item []
\begin{eqnarray*}
y-y_{1} & = & m(x-x_{1})\\
y-4 & = & -2(x-13)\\
y-4 & = & -2x+26\\
y & = & -2x+30
\end{eqnarray*}
\vfill
\end{enumerate}
\item Passing through $(-3,0)$ and $(-6,4)$.

\begin{enumerate}[resume]
\item []$m=\frac{0-4}{-3-(-6)}=\frac{-4}{3}=-\frac{4}{3}$
\begin{eqnarray*}
y-y_{1} & = & m(x-x_{1})\\
y-0 & = & -\frac{4}{3}(x-(-3))\\
y & = & -\frac{4}{3}x+3
\end{eqnarray*}
\vfill
\end{enumerate}
\end{enumerate}
\textbf{Find the equation of the least squares line. }
\begin{enumerate}[resume]
\item \label{enu:The-total-amount}The paired data below consists of the
number or Florida pleasure boats (in tens of thousands) and the number
of yearly watercraft-related manatee deaths from 1995-2000.

\begin{enumerate}[resume]
\item []%
\begin{tabular}{c|c|c|c|c|c|c}
Hours $(x)$ & 68 & 67 & 70 & 73 & 81 & 84\tabularnewline
\hline 
Score $(y)$ & 53 & 35 & 49 & 60 & 67 & 78\tabularnewline
\end{tabular}

\begin{enumerate}[resume]
\item [TI-83/84:] 

\begin{enumerate}[resume]
\item [Step 1:]STATS>EDIT enter the ``Hours'' into L1 and ``Score''
into L2. 
\item [Step 2:]STAT>CALC>LinReg \emph{or }STAT>TESTS>LinRegTTest>Caculate
\item []$a\approx1.9415$ and $b\approx-86.3495\Rightarrow y=1.9415x-86.35$\vfill
\end{enumerate}
\end{enumerate}
\end{enumerate}
\end{enumerate}
\textbf{Solve the system of two equations in two unknowns.}
\begin{enumerate}[resume]
\item $\begin{array}{r}
x+y=5\\
x-y=10
\end{array}$

\begin{enumerate}[resume]
\item []$\begin{array}{r}
x+y=5\\
x-y=10
\end{array}\Rightarrow\begin{array}{r}
x+y=5\\
2x+0=15
\end{array}\Rightarrow2x=15\Rightarrow x=\frac{15}{2}$\\
Substituting $x$: $\left(\frac{15}{2}\right)+y=5\Rightarrow y=5-\frac{15}{2}\Rightarrow y=-\frac{5}{2}$
\item []Solution: $(\frac{15}{2},-\frac{5}{2})$
\end{enumerate}
\item $\begin{array}{r}
x+y=2\\
2x+2y=5
\end{array}$

\begin{enumerate}[resume]
\item []$\begin{array}{r}
x+y=2\\
2x+2y=5
\end{array}\Rightarrow\begin{array}{r}
2x+2y=4\\
2x+2y=5
\end{array}\Rightarrow\begin{array}{r}
2x+2y=4\\
0+0=1
\end{array}\Rightarrow0=1\Rightarrow\mbox{No solution exists}$\vfill
\end{enumerate}
\end{enumerate}
\textbf{Solve the problem.}
\begin{enumerate}[resume]
\item There are 1000 people at a basketball game. The admission prices was
\$20 for non-students and \$10 for students. The admission receipts
were \$17,500. How many students and non-students attended the game?

\begin{enumerate}[resume]
\item []$20x+20y=20000\begin{array}{r}
x+y=1000\\
20x+10y=17500
\end{array}\Rightarrow\begin{array}{r}
0+10y=2500\\
20x+10y=17500
\end{array}\Rightarrow10y=2500$$\Rightarrow y=25020x+10y=17500$\\
$y=250\Rightarrow x+(250)=1000\Rightarrow x=750$. \\
So there were $750$ non-students and $250$ students at the game.\vfill
\end{enumerate}
\item Two pounds of rice and three pounds of potatoes cost a total of \$8,
and three pounds of rice and one pound of potatoes cost a total of
\$5. What is the price per pound for rice and the price per pound
of potatoes?

\begin{enumerate}[resume]
\item []$\begin{array}{r}
2x+3y=8\\
3x+1y=5
\end{array}\Rightarrow\begin{array}{r}
2x+3y=8\\
9x+3y=15
\end{array}\Rightarrow\begin{array}{r}
2x+3y=8\\
7x+0y=7
\end{array}\Rightarrow x=1\Rightarrow3(1)+y=5\Rightarrow y=2$\\
So rice costs \$1 per pound and potatoes cost \$2 per pound. \vfill
\end{enumerate}
\item A car manufacturer has decided to come out with a new type of car.
The fixed cost for the production of this new car is \$250,000 and
the variable costs are \$15,000 per car. The company expects to sell
the cars for \$20,000 each. Formulate a function $C(x)$ for the total
cost of producing $x$ new cars and a function $R(x)$ for the total
revenue generated by the sales of $x$ cars. 

\begin{enumerate}[resume]
\item []$C(x)=15000x+250000$.
\item []$R(x)=20000x$\vfill
\end{enumerate}
\item After 5 years as an actuary, Sue's salary was \$100,000. After 10
years her salary had risen to \$200,000. Let $y$ represent her salary
after $x$ years on the job. Assuming that a change in her salary
over time can be approximated by a straight line, give an equation
for this line in the form $y=mx+b$.

\begin{enumerate}[resume]
\item []Years and salary give us two points: $(5,100000)$ and $(10,200000)$.\\
$m=\frac{200000-100000}{10-5}=\frac{100000}{5}=20000$. \\
\begin{eqnarray*}
y-100000 & = & 20000(x-5)\\
y-100000 & = & 20000x-100000\\
y & = & 20000x
\end{eqnarray*}
\vfill
\end{enumerate}
\item A package delivery company delivers packages at a cost of \$5.00 per
package. The fixed cost to run the delivery truck is \$450 per day.
If the company charges \$7.00 per package, how many packages must
be delivered to make a profit of \$500?

\begin{enumerate}[resume]
\item []
\begin{eqnarray*}
P(x) & = & R(x)-C(x)\\
 & = & \left[7x\right]-\left[5x+450\right]\\
 & = & 7x-5x-450\\
 & = & 2x-450
\end{eqnarray*}
need to solve $P(x)=500$,
\begin{eqnarray*}
500 & = & 2x-450\\
950 & = & 2x\\
475 & = & x
\end{eqnarray*}
so $475$ packages.\vfill
\end{enumerate}
\item A package delivery company delivers packages at a cost of \$5 per
package . The fixed cost to run the delivery truck is \$450 per day.
If the company charges \$7 per package, how many packages must be
delivered to break even?

\begin{enumerate}[resume]
\item []The profit function is the same as above. So we need to solve $P(x)=0$.
\begin{eqnarray*}
0 & = & 2x-450\\
450 & = & 2x\\
225 & = & x
\end{eqnarray*}
so $225$ packages.\vfill
\end{enumerate}
\item Let the supply and demand functions for gasoline be given by 
\[
p=S(q)=\frac{4}{5}q\mbox{ and }p=D(q)=300-\frac{2}{5}q
\]
where $p$ is the price in dollars and $q$ is the number of 42-gallon
barrels. Find the equilibrium quantity and the equilibrium price.

\begin{enumerate}[resume]
\item []
\begin{eqnarray*}
S(q) & = & D(q)\\
\frac{4}{5}q & = & 300-\frac{2}{5}q\\
\frac{6}{5}q & = & 300\\
q & = & \frac{5}{6}300=250
\end{eqnarray*}
Equilibrium quantity: $q=250$
\[
\frac{4}{5}(250)=200
\]
Equilibrium price: $p=\$200$\vfill
\end{enumerate}
\end{enumerate}
\textbf{The sizes of two matrices $A$ and $B$ are given. Find the
sizes of the product $AB$ and the product $BA$, whenever these products
exist.}
\begin{enumerate}[resume]
\item $A$ is $2\times5$, and $B$ is $5\times3$.

\begin{enumerate}[resume]
\item $AB$:

\begin{enumerate}[resume]
\item []$2\times3$
\end{enumerate}
\item $BA$:

\begin{enumerate}[resume]
\item []Does not exist\vfill
\end{enumerate}
\end{enumerate}
\end{enumerate}
\textbf{Find the matrix if possible.}
\begin{enumerate}[resume]
\item $\begin{bmatrix}\begin{array}{rr}
3 & 0\\
-2 & 1\\
0 & 3
\end{array}\end{bmatrix}\begin{bmatrix}\begin{array}{rrr}
1 & 0 & 3\\
0 & 1 & 0
\end{array}\end{bmatrix}$

\begin{enumerate}[resume]
\item []$\begin{bmatrix}\begin{array}{rrr}
3(1)+0(0) & 3(0)+0(1) & 3(3)+0(0)\\
-2(1)+1(0) & -2(0)+1(1) & -2(3)+1(0)\\
0(1)+3(0) & 0(0)+3(1) & 0(3)+3(0)
\end{array}\end{bmatrix}=\begin{bmatrix}\begin{array}{rrr}
3 & 0 & 9\\
-2 & 1 & -6\\
0 & 3 & 0
\end{array}\end{bmatrix}$
\end{enumerate}
\item $\begin{bmatrix}\begin{array}{rr}
3 & 0\\
-2 & 1\\
0 & 3
\end{array}\end{bmatrix}\begin{bmatrix}\begin{array}{rr}
1 & 0\\
0 & 0\\
0 & 0
\end{array}\end{bmatrix}$

\begin{enumerate}[resume]
\item Not possible. You can not multiply a $(3\times2)$ and $(3\times2)$.
\vfill
\end{enumerate}
\item $\begin{bmatrix}\begin{array}{rrr}
a & b & 1\\
0 & c & 0\\
0 & 1 & d
\end{array}\end{bmatrix}\begin{bmatrix}\begin{array}{rrr}
0 & 0 & 0\\
0 & 0 & 0\\
0 & 0 & 0
\end{array}\end{bmatrix}$ where $a,b,c,d$ are real numbers.

\begin{enumerate}[resume]
\item []$\begin{bmatrix}\begin{array}{rrr}
a & b & 1\\
0 & c & 0\\
0 & 1 & d
\end{array}\end{bmatrix}\begin{bmatrix}\begin{array}{rrr}
0 & 0 & 0\\
0 & 0 & 0\\
0 & 0 & 0
\end{array}\end{bmatrix}=\begin{bmatrix}\begin{array}{rrr}
0 & 0 & 0\\
0 & 0 & 0\\
0 & 0 & 0
\end{array}\end{bmatrix}$\vfill
\end{enumerate}
\end{enumerate}
\textbf{Solve the problem.}
\begin{enumerate}[resume]
\item Let $A=\begin{bmatrix}\begin{array}{rr}
-1 & 3\\
2 & 0
\end{array}\end{bmatrix}$ Find $-2A$. 

\begin{enumerate}[resume]
\item []$-2A=\begin{bmatrix}\begin{array}{rr}
-1(-2) & 3(-2)\\
2(-2) & 0(-2)
\end{array}\end{bmatrix}=\begin{bmatrix}\begin{array}{rr}
2 & -6\\
-4 & -0
\end{array}\end{bmatrix}$\vfill
\end{enumerate}
\item Let $A=\begin{bmatrix}\begin{array}{rr}
-1 & 3\\
2 & 0
\end{array}\end{bmatrix}$ and $B=\begin{bmatrix}\begin{array}{rr}
0 & 2\\
5 & -1
\end{array}\end{bmatrix}$. Find $A+B$. 

\begin{enumerate}[resume]
\item []$A+B=\begin{bmatrix}\begin{array}{rr}
-1+0 & 3+2\\
2+5 & 0-1
\end{array}\end{bmatrix}=\begin{bmatrix}\begin{array}{rr}
-1 & 5\\
7 & -1
\end{array}\end{bmatrix}$\vfill
\end{enumerate}
\item Let $A=\begin{bmatrix}\begin{array}{rr}
-1 & 3\\
2 & 0
\end{array}\end{bmatrix}$ and $B=\begin{bmatrix}\begin{array}{rr}
0 & 2\\
5 & -1
\end{array}\end{bmatrix}$. Find $2A+B$. 

\begin{enumerate}[resume]
\item []$2A+B=\begin{bmatrix}\begin{array}{rr}
-2 & 6\\
4 & 0
\end{array}\end{bmatrix}+\begin{bmatrix}\begin{array}{rr}
0 & 2\\
5 & -1
\end{array}\end{bmatrix}=\begin{bmatrix}\begin{array}{rr}
-2+0 & 6+2\\
4+5 & 0-1
\end{array}\end{bmatrix}=\begin{bmatrix}\begin{array}{rr}
-2 & 8\\
9 & -1
\end{array}\end{bmatrix}$\vfill
\end{enumerate}
\end{enumerate}
\textbf{Bonus}
\begin{enumerate}[resume]
\item [$\bigstar$]Find the inverse ($A^{-1}$) of $A=\begin{bmatrix}\begin{array}{rr}
1 & 0\\
2 & 2
\end{array}\end{bmatrix}$ without using a calculator/computer.

\begin{enumerate}[resume]
\item []$\begin{bmatrix}\begin{array}{rr|rr}
1 & 0 & 1 & 0\\
2 & 2 & 0 & 1
\end{array}\end{bmatrix}\rightsquigarrow\begin{bmatrix}\begin{array}{rr|rr}
1 & 0 & 1 & 0\\
0 & 2 & -2 & 1
\end{array}\end{bmatrix}\begin{array}{l}
\mbox{}\\
\mbox{R2}-2\mbox{R1}
\end{array}\rightsquigarrow\begin{bmatrix}\begin{array}{rr|rr}
1 & 0 & 1 & 0\\
0 & 1 & -1 & -\nicefrac{1}{2}
\end{array}\end{bmatrix}\begin{array}{l}
\mbox{}\\
\nicefrac{1}{2}\mbox{R2}
\end{array}$\\
So $A^{-1}=\begin{bmatrix}\begin{array}{rr}
1 & 0\\
-1 & -\nicefrac{1}{2}
\end{array}\end{bmatrix}$\vfill
\end{enumerate}
\end{enumerate}

\end{document}
